\clearpage\pagestyle{otMain}\setlength{\otBlocEspacement}{2.5mm}\fontsize{11}{15}\selectfont\corrigesChapitre{7}{Proportionnalité et pourcentages}
\begin{blocLivre}{cor-7-1}\begin{corrige}{1}{Passer à l’unité}\demarche{Calculer le prix d’une unité.}
Une bouteille coûte $9/6=1{,}50$ euro. Dix bouteilles coûtent $10\times1{,}50=\res{15}$ euros. On suppose le prix unitaire constant.
\end{corrige}\end{blocLivre}
\begin{blocLivre}{cor-7-2}\begin{corrige}{2}{Compléter un tableau}\demarche{Préciser le sens du coefficient.}
Le coefficient masse vers prix est $7/2=3{,}5$ euros par kg. Les prix sont donc $5\times3{,}5=\res{17{,}50}$ euros et $8\times3{,}5=\res{28}$ euros.
\end{corrige}\end{blocLivre}
\begin{blocLivre}{cor-7-3}\begin{corrige}{3}{Une part en pourcentage}\demarche{Multiplier la quantité par la fraction correspondante.}
$18\,\%=18/100$. La part cherchée vaut $(18/100)\times250=\res{45}$.
\end{corrige}\end{blocLivre}
\begin{blocLivre}{cor-7-4}\begin{corrige}{4}{Retrouver le pourcentage}\demarche{Diviser la partie par le total.}
Le rapport partie sur total est $12/32=0{,}375$. Il correspond à $\res{37{,}5\,\%}$. Le dénominateur est l’effectif total, pas celui des autres élèves.
\end{corrige}\end{blocLivre}
\begin{blocLivre}{cor-7-5}\begin{corrige}{5}{Lire une échelle}\demarche{Faire d’abord le rapport avec des unités identiques.}
La longueur réelle est $3{,}6\times500=1800$ cm, soit $\res{18\ \mathrm m}$. La conversion en mètres vient après le calcul dans la même unité.
\end{corrige}\end{blocLivre}
\begin{blocLivre}{cor-7-6}\begin{corrige}{6}{Une quatrième proportionnelle}\demarche{Comparer passage à l’unité et calcul direct.}
Par unité : $14{,}4/8=1{,}8$ kg, puis $5\times1{,}8=\res9$ kg. Directement, on calcule $5\times14{,}4/8=9$ kg. Les deux méthodes utilisent le même prix de passage, ici une masse unitaire.
\end{corrige}\end{blocLivre}
\begin{blocLivre}{cor-7-7}\begin{corrige}{7}{Une remise}\demarche{Calculer la baisse avant le prix final.}
La remise vaut $80\times12/100=9{,}60$ euros. On la retire : $80-9{,}60=\res{70{,}40}$ euros. Une remise est une diminution.
\end{corrige}\end{blocLivre}
\begin{blocLivre}{cor-7-8}\begin{corrige}{8}{Un tarif non proportionnel}\demarche{Tester un doublement.}
Pour 1 h : $4+3=\res7$ euros. Pour 2 h : $4+6=\res{10}$ euros. Doubler le temps ne double pas 7 euros, donc il n’y a pas proportionnalité. Les frais fixes l’empêchent.
\end{corrige}\end{blocLivre}
\begin{blocLivre}{cor-7-9}\begin{corrige}{9}{Un graphique}\demarche{Vérifier l’alignement avec l’origine puis lire les axes.}
Les points sont sur une droite passant par l’origine, ce qui représente une proportionnalité dans ce modèle. À 3 kg, l’ordonnée est $\res{7{,}50}$ euros. Le coefficient est 2,50 euros par kg.
\label{fig-cor-7-9}\figureGeo[corrige]{proportion-prix}
\end{corrige}\end{blocLivre}
\begin{blocLivre}{cor-7-10}\begin{corrige}{10}{Une formule en sciences}\demarche{Isoler la grandeur cherchée dans la formule.}
$U=15\times0{,}4=\res6$ V. Pour 9 V, $9=15I$, donc $I=9/15=\res{0{,}6}$ A. Le coefficient tension sur intensité est constant et vaut 15.
\end{corrige}\end{blocLivre}
\begin{blocLivre}{cor-7-11}\begin{corrige}{11}{Deux remises successives}\demarche{Changer de base pour la seconde remise.}
Après 20\,\%, le prix est 80 euros. La seconde baisse vaut $10\,\%$ de 80, soit 8 euros : le prix final est $\res{72}$ euros. La baisse totale vaut 28 euros sur 100, donc $\res{28\,\%}$, et non 30\,\%.
\end{corrige}\end{blocLivre}
\begin{blocLivre}{cor-7-12}\begin{corrige}{12}{Une recette et des lots}\demarche{Calculer la quantité puis arrondir le nombre de lots vers le haut.}
Par personne, il faut $450/6=75$ g. Pour 14 personnes, $75\times14=1050$ g. Deux paquets ne donnent que 1000 g : il faut $\res3$ paquets. La quantité achetée peut dépasser la quantité nécessaire.
\end{corrige}\end{blocLivre}
\begin{blocLivre}{cor-7-13}\begin{corrige}{13}{Un pourcentage inverse}\demarche{Identifier la part conservée du prix initial.}
Le prix payé représente 75\,\% du prix initial $P$. On écrit $(75/100)P=54$, d’où $P=54\times100/75=\res{72}$ euros. Vérification : 25\,\% de 72 vaut 18, et $72-18=54$.
\end{corrige}\end{blocLivre}
\begin{blocLivre}{cor-7-14}\begin{corrige}{14}{Un arc de cercle}\demarche{Prendre la fraction d’un tour complet.}
La circonférence vaut $2\pi\times6=12\pi$ cm. L’arc en représente $75/360$, donc sa longueur vaut $12\pi\times75/360=\res{5\pi/2\ \mathrm{cm}}$, soit environ $\res{7{,}9\ \mathrm{cm}}$.
\label{fig-cor-7-14}\figureGeo[corrige]{arc-75}
\end{corrige}\end{blocLivre}
\begin{blocLivre}{cor-7-15}\begin{corrige}{15}{Deux groupes réunis}\demarche{Revenir aux effectifs avant de réunir les groupes.}
Il y a $20\times0{,}4=8$ filles et $30\times0{,}6=18$ femmes, soit 26 personnes sur 50. Le pourcentage est $26/50\times100=\res{52\,\%}$. La moyenne simple de 40\,\% et 60\,\% ignorerait les effectifs différents.
\end{corrige}\end{blocLivre}
